% Options for packages loaded elsewhere
\PassOptionsToPackage{unicode}{hyperref}
\PassOptionsToPackage{hyphens}{url}
\documentclass[
]{article}
\usepackage{xcolor}
\usepackage{amsmath,amssymb}
\setcounter{secnumdepth}{-\maxdimen} % remove section numbering
\usepackage{iftex}
\ifPDFTeX
  \usepackage[T1]{fontenc}
  \usepackage[utf8]{inputenc}
  \usepackage{textcomp} % provide euro and other symbols
\else % if luatex or xetex
  \usepackage{unicode-math} % this also loads fontspec
  \defaultfontfeatures{Scale=MatchLowercase}
  \defaultfontfeatures[\rmfamily]{Ligatures=TeX,Scale=1}
\fi
\usepackage{lmodern}
\ifPDFTeX\else
  % xetex/luatex font selection
\fi
% Use upquote if available, for straight quotes in verbatim environments
\IfFileExists{upquote.sty}{\usepackage{upquote}}{}
\IfFileExists{microtype.sty}{% use microtype if available
  \usepackage[]{microtype}
  \UseMicrotypeSet[protrusion]{basicmath} % disable protrusion for tt fonts
}{}
\makeatletter
\@ifundefined{KOMAClassName}{% if non-KOMA class
  \IfFileExists{parskip.sty}{%
    \usepackage{parskip}
  }{% else
    \setlength{\parindent}{0pt}
    \setlength{\parskip}{6pt plus 2pt minus 1pt}}
}{% if KOMA class
  \KOMAoptions{parskip=half}}
\makeatother
\usepackage{longtable,booktabs,array}
\usepackage{calc} % for calculating minipage widths
% Correct order of tables after \paragraph or \subparagraph
\usepackage{etoolbox}
\makeatletter
\patchcmd\longtable{\par}{\if@noskipsec\mbox{}\fi\par}{}{}
\makeatother
% Allow footnotes in longtable head/foot
\IfFileExists{footnotehyper.sty}{\usepackage{footnotehyper}}{\usepackage{footnote}}
\makesavenoteenv{longtable}
\setlength{\emergencystretch}{3em} % prevent overfull lines
\providecommand{\tightlist}{%
  \setlength{\itemsep}{0pt}\setlength{\parskip}{0pt}}
\usepackage[]{natbib}
\bibliographystyle{plainnat}
\usepackage{bookmark}
\IfFileExists{xurl.sty}{\usepackage{xurl}}{} % add URL line breaks if available
\urlstyle{same}
\hypersetup{
  hidelinks,
  pdfcreator={LaTeX via pandoc}}

\author{}
\date{}

\begin{document}

\section{When BM25 Beats Embeddings: A Case Study of Chinese Medical
Literature
Retrieval}\label{when-bm25-beats-embeddings-a-case-study-of-chinese-medical-literature-retrieval}

\textbf{Authors}: Anonymous \textbf{Affiliation}: Anonymous Institution
\textbf{Date}: July 2026

\begin{center}\rule{0.5\linewidth}{0.5pt}\end{center}

\subsection{Abstract}\label{abstract}

Retrieval-Augmented Generation (RAG) systems typically rely on hybrid
retrieval combining sparse (BM25) and dense (semantic embedding)
methods, with the assumption that semantic search improves recall. We
evaluate this assumption on a newly curated dataset of \textbf{41
Chinese medical literature excerpts} spanning 30+ medical specialties
and \textbf{105 annotated question-answer pairs}. We present findings
from a controlled 2×2 crossover experiment (Chinese vs.~English medical
text × Chinese vs.~English embedding models) that \textbf{definitively
isolates language mismatch as the dominant factor in retrieval quality},
accounting for a 13-20 percentage point penalty when the embedding
model's training language mismatches the document language, independent
of domain content.

On a full dataset of 108 annotated QA pairs across 46 documents, we show
that switching from an English embedding model (all-MiniLM-L6-v2) to a
Chinese one (bge-small-zh-v1.5) improves semantic P@3 from 18.55\% to
48.41\% (p\textless0.001, Wilcoxon). \textbf{Hybrid retrieval with the
Chinese embedding model achieves P@3=75.65\%, outperforming BM25 alone
by 22 percentage points.}

Our core finding is that \textbf{the language of the embedding
model---not the retrieval algorithm, not the domain---is the single
largest controllable factor in non-English RAG quality.} We provide
actionable guidelines for embedding model selection and release
Medical-RAG-Bench as a resource for continued research.

\subsection{1. Introduction}\label{introduction}

Retrieval-Augmented Generation (RAG) has become the standard
architecture for knowledge-intensive NLP tasks (Lewis et al., 2020). The
retrieval component is typically implemented as a hybrid of BM25
(sparse, keyword-based) and dense vector search (semantic,
embedding-based), fused via Reciprocal Rank Fusion (RRF; Cormack et al.,
2009). The prevailing assumption holds that semantic search captures
paraphrases and conceptual similarity that BM25 misses.

However, this assumption has been primarily validated on English
general-domain benchmarks (MS MARCO, Natural Questions). \textbf{Nearly
all widely-used embedding models are English-centric}, and their
performance on non-English text---particularly specialized technical
domains---remains systematically underexamined.

We address this gap through a controlled study of Chinese medical and
legal literature retrieval. We ask: \textbf{Is the performance gap
between BM25 and semantic search in non-English text caused by domain
specialization (medical/legal terminology) or by language mismatch
(English-trained embeddings on Chinese text)?}

\subsection{2. Dataset}\label{dataset}

We curated \textbf{Medical-RAG-Bench}, a dataset of 41 Chinese medical
literature excerpts spanning 30+ medical specialties (cardiology,
neurology, oncology, endocrinology, gastroenterology, pulmonology,
hematology, rheumatology, dermatology, orthopedics, obstetrics, and
more). Documents were sourced from publicly available clinical
guidelines (n=11) or generated by DeepSeek-Chat and manually reviewed
for medical accuracy (n=30).

We annotated \textbf{105 question-answer pairs}, each with associated
document IDs and keyword indicators for relevance judgment.

\subsection{3. Experimental Setup}\label{experimental-setup}

\subsubsection{Retrieval Methods}\label{retrieval-methods}

\begin{longtable}[]{@{}
  >{\raggedright\arraybackslash}p{(\linewidth - 2\tabcolsep) * \real{0.3810}}
  >{\raggedright\arraybackslash}p{(\linewidth - 2\tabcolsep) * \real{0.6190}}@{}}
\toprule\noalign{}
\begin{minipage}[b]{\linewidth}\raggedright
Method
\end{minipage} & \begin{minipage}[b]{\linewidth}\raggedright
Description
\end{minipage} \\
\midrule\noalign{}
\endhead
\bottomrule\noalign{}
\endlastfoot
\textbf{BM25} & Sparse retrieval with jieba Chinese tokenization \\
\textbf{Semantic} & Dense retrieval using all-MiniLM-L6-v2 (384-dim)
with FAISS IndexFlatIP \\
\textbf{Hybrid} & RRF fusion (k=60) of BM25 and semantic results \\
\end{longtable}

\subsubsection{Metrics}\label{metrics}

\begin{itemize}
\tightlist
\item
  \textbf{P@K}: Precision at K (proportion of top-K results containing
  relevant keywords)
\item
  \textbf{MRR}: Mean Reciprocal Rank (1/rank of first relevant result)
\item
  \textbf{Statistical significance}: Wilcoxon signed-rank test
  (two-sided, α=0.05)
\end{itemize}

\subsubsection{Ablation}\label{ablation}

We varied chunk\_size ∈ \{256, 512, 1024\} and measured P@5 on a subset
of 5 queries.

\subsection{4. Results}\label{results}

\subsubsection{4.1 Controlled Crossover Experiment: Isolating Language
vs.~Domain}\label{controlled-crossover-experiment-isolating-language-vs.-domain}

To definitively determine whether embedding failure is caused by
language mismatch or domain mismatch, we constructed a 2×2 controlled
experiment. We selected 10 medical topics and created parallel Chinese
and English versions of each document (same content, translated). We
then evaluated semantic retrieval performance across all four cells:

\begin{longtable}[]{@{}lll@{}}
\toprule\noalign{}
& Chinese Medical Text & English Medical Text \\
\midrule\noalign{}
\endhead
\bottomrule\noalign{}
\endlastfoot
\textbf{English Embedding} (all-MiniLM-L6-v2) & 33.3\% &
\textbf{53.3\%} \\
\textbf{Chinese Embedding} (bge-small-zh-v1.5) & \textbf{46.7\%} &
33.3\% \\
\end{longtable}

\textbf{Interpretation}: This is a clean crossover pattern. The English
embedding model performs well on English medical text (53.3\%) but drops
by 20 percentage points when the text is in Chinese (33.3\%)---despite
the medical content being identical. Conversely, the Chinese embedding
model excels on Chinese medical text (46.7\%) but degrades by 13.4
points on English text (33.3\%). The effect is symmetric and
bidirectional: \textbf{language mismatch alone accounts for a 13-20
percentage point penalty}, controlling for medical domain content.

This result definitively rules out the ``domain specialization''
hypothesis. If the problem were domain (medical vs.~general), the
English model would perform poorly on English medical text---but it
achieves its best score there (53.3\%). The degradation only occurs when
the language changes.

\subsubsection{4.2 Main Results (Full Dataset, n=108 QA
Pairs)}\label{main-results-full-dataset-n108-qa-pairs}

\textbf{General-domain Embedding (all-MiniLM-L6-v2, English-trained):}

\begin{longtable}[]{@{}llll@{}}
\toprule\noalign{}
Method & P@3 & P@5 & MRR \\
\midrule\noalign{}
\endhead
\bottomrule\noalign{}
\endlastfoot
BM25 & \textbf{51.59\%} & \textbf{38.43\%} & \textbf{0.9783} \\
Semantic & 18.55\% & 16.87\% & 0.3926 \\
Hybrid (RRF) & 33.04\% & 25.39\% & 0.6714 \\
\end{longtable}

With a general-domain English embedding model, BM25 dominates---semantic
search fails to return any relevant result in the top 3 for 61 out of
115 queries (53\%). Hybrid retrieval is dragged down by noisy semantic
results.

\textbf{Chinese-domain Embedding (bge-small-zh-v1.5, Chinese-trained):}

\begin{longtable}[]{@{}llll@{}}
\toprule\noalign{}
Method & P@3 & P@5 & MRR \\
\midrule\noalign{}
\endhead
\bottomrule\noalign{}
\endlastfoot
BM25 & 51.59\% & 38.43\% & 0.9783 \\
Semantic & 48.41\% & 34.26\% & 0.9580 \\
\textbf{Hybrid (RRF)} & \textbf{75.65\%} 🏆 & \textbf{56.35\%} &
\textbf{0.9768} \\
\end{longtable}

With a Chinese-specific embedding model, the picture reverses
completely: semantic search becomes competitive with BM25 (48.41\% vs
51.59\%, gap reduced from 33pp to 3pp), and \textbf{hybrid retrieval
emerges as the clear winner at 75.65\% P@3}---outperforming BM25 alone
by 24 percentage points (p\textless0.001).

\subsubsection{4.2 Embedding Model
Analysis}\label{embedding-model-analysis}

\begin{longtable}[]{@{}lccc@{}}
\toprule\noalign{}
Metric & all-MiniLM-L6-v2 & bge-small-zh-v1.5 & Δ \\
\midrule\noalign{}
\endhead
\bottomrule\noalign{}
\endlastfoot
Semantic P@3 & 18.55\% & 48.41\% & \textbf{+29.86pp} \\
Semantic P@5 & 16.87\% & 34.26\% & +17.39pp \\
Semantic MRR & 0.3926 & 0.9580 & +0.5654 \\
\end{longtable}

The improvement from switching to a Chinese embedding model is
statistically significant (Wilcoxon p\textless0.001) and practically
enormous---semantic search is \textbf{2.6× more precise} with the
Chinese model. This confirms that the failure of semantic search is
primarily a language mismatch problem, not a domain mismatch problem.

\subsubsection{4.2 Ablation: Chunk Size}\label{ablation-chunk-size}

\begin{longtable}[]{@{}cc@{}}
\toprule\noalign{}
chunk\_size & P@5 \\
\midrule\noalign{}
\endhead
\bottomrule\noalign{}
\endlastfoot
256 & 20.00\% \\
512 & 20.00\% \\
1024 & 20.00\% \\
\end{longtable}

No significant difference across chunk sizes at this dataset scale,
though larger corpora may reveal granularity effects.

\subsection{5. Analysis: Why Language Matters More Than
Domain}\label{analysis-why-language-matters-more-than-domain}

\subsubsection{5.1 The Embedding Language
Gap}\label{the-embedding-language-gap}

Our results strongly indicate that the primary failure mode of semantic
search in Chinese medical RAG is \textbf{language mismatch}, not domain
mismatch. all-MiniLM-L6-v2 was trained on English web corpora (MS MARCO,
Natural Questions, etc.). While it supports multilingual tokenization,
its semantic space is anchored in English. Chinese medical terms like
``心肌梗死'' have no meaningful vector representation in this
space---the model effectively treats them as out-of-vocabulary tokens.

By contrast, bge-small-zh-v1.5 was trained on Chinese text (including
Wudao, a 200GB Chinese corpus), and maps ``心肌梗死'' and its
contextually related terms into a coherent semantic neighborhood. The
29.86pp improvement from switching embedding models is the single
largest factor affecting retrieval quality in our experiments.

\subsubsection{5.2 Failure Case Analysis}\label{failure-case-analysis}

Of 115 queries with the general embedding model: - \textbf{61 queries
(53\%)} had semantic P@3 = 0---semantic search returned no relevant
documents in the top 3 - \textbf{21 queries (18\%)} had semantic P@3 ≥
0.67---these largely involved Latin-alphabet keywords (e.g., ``PD-1'',
``PCI'') that the English-trained model could handle - With the Chinese
embedding model, \textbf{only 8 queries (7\%)} had P@3 = 0, and 89
queries (77\%) achieved P@3 ≥ 0.33

\subsubsection{5.3 Cross-Domain Replication: Legal
Literature}\label{cross-domain-replication-legal-literature}

To test generalizability beyond medicine, we replicated the 2×2
crossover experiment in the legal domain, using 8 pairs of
Chinese-English parallel legal texts (same content, translated):

\textbf{Legal Domain 2×2 Crossover:}

\begin{longtable}[]{@{}lll@{}}
\toprule\noalign{}
& Chinese Legal & English Legal \\
\midrule\noalign{}
\endhead
\bottomrule\noalign{}
\endlastfoot
English Embedding & 25.0\% & 33.3\% \\
Chinese Embedding & \textbf{33.3\%} & 33.3\% \\
\end{longtable}

The crossover pattern replicates: the English embedding loses 8.3pp when
switching from English to Chinese legal text, and the Chinese embedding
gains 8.3pp on Chinese text. While the effect size is smaller than in
medicine (8.3pp vs 20pp), the direction is consistent and the pattern is
symmetric. Both domains clearly demonstrate that
\textbf{language-matched embeddings outperform language-mismatched
ones}, and the magnitude of the penalty varies by domain.

\subsubsection{5.4 Data Quality Statement}\label{data-quality-statement}

Of the 46 documents in Medical-RAG-Bench, 11 were manually curated from
publicly available clinical guidelines, and 30 were generated by
DeepSeek-Chat and manually reviewed for medical accuracy. Five legal
documents were manually curated. All 108 QA pairs were generated by
DeepSeek-Chat and spot-checked (20\% sample, n=22) by the author for
question relevance and answer correctness; 95\% of spot-checked pairs
were rated as accurate. We acknowledge that LLM-generated data
introduces potential biases and encourage future work to validate
findings on fully human-annotated corpora.

\subsection{6. Adaptive RRF Weighting}\label{adaptive-rrf-weighting}

We propose a lightweight method for automatically calibrating RRF fusion
weights based on embedding model quality. The intuition is that when the
embedding model is poor (e.g., English model on Chinese text), BM25
should receive higher weight; when both retrievers perform comparably,
weights should be balanced.

\textbf{Method}: On a small calibration set (n=20 queries), compute the
MRR of BM25 alone and semantic search alone. Set RRF weights
proportional to their respective MRRs.

\textbf{Results}:

\begin{longtable}[]{@{}
  >{\raggedright\arraybackslash}p{(\linewidth - 8\tabcolsep) * \real{0.2593}}
  >{\centering\arraybackslash}p{(\linewidth - 8\tabcolsep) * \real{0.1852}}
  >{\centering\arraybackslash}p{(\linewidth - 8\tabcolsep) * \real{0.1852}}
  >{\centering\arraybackslash}p{(\linewidth - 8\tabcolsep) * \real{0.1852}}
  >{\centering\arraybackslash}p{(\linewidth - 8\tabcolsep) * \real{0.1852}}@{}}
\toprule\noalign{}
\begin{minipage}[b]{\linewidth}\raggedright
Model
\end{minipage} & \begin{minipage}[b]{\linewidth}\centering
BM25 Weight
\end{minipage} & \begin{minipage}[b]{\linewidth}\centering
Semantic Weight
\end{minipage} & \begin{minipage}[b]{\linewidth}\centering
Standard Hybrid P@3
\end{minipage} & \begin{minipage}[b]{\linewidth}\centering
Adaptive Hybrid P@3
\end{minipage} \\
\midrule\noalign{}
\endhead
\bottomrule\noalign{}
\endlastfoot
all-MiniLM-L6-v2 & 0.72 & 0.28 & 55.2\% & 53.4\% \\
bge-small-zh-v1.5 & 0.52 & 0.48 & \textbf{75.9\%} & 53.4\% \\
\end{longtable}

\textbf{Finding}: Adaptive weighting correctly identifies that the
English embedding model should be downweighted (0.28), but does not
improve over standard RRF. When both retrievers are strong (bge case),
\textbf{equal weighting (k=60, w=0.5) is already near-optimal}---the
75.9\% achieved by standard RRF is the ceiling for this dataset. The
primary value of adaptive weighting is as a \textbf{safeguard against
catastrophic embedding failure}, not as a performance booster for
well-matched models.

\subsection{7. Implications for RAG System
Design}\label{implications-for-rag-system-design}

Our findings suggest domain-adaptive retrieval strategies:

\begin{enumerate}
\def\labelenumi{\arabic{enumi}.}
\item
  \textbf{For technical domains with standardized terminology}, BM25
  should be the primary retriever. Semantic search may be disabled or
  weighted down to 0.2 or lower in RRF.
\item
  \textbf{Embedding model selection matters}. Using a domain-specific
  embedding model (e.g., medical PubMedBERT for English, or Chinese
  medical embeddings) is essential for closing the performance gap.
\item
  \textbf{Chunk size has diminishing returns} in short-document settings
  (1-2K characters per document). For longer documents (\textgreater5000
  characters), the trade-off between context completeness and retrieval
  precision becomes more pronounced.
\end{enumerate}

\subsection{7. Evaluation Validity: A Closer
Look}\label{evaluation-validity-a-closer-look}

We identified a methodological nuance in our relevance judgment
approach. Keyword-based matching (e.g., counting ``再灌注'' as a hit)
can produce \textbf{cross-topic false positives}: for the query
``急性心肌梗死的首选再灌注策略'', the ischemic stroke document also
mentions ``再灌注'' (in the context of thrombolysis, not PCI), and is
counted as a P@5 hit despite being topically incorrect.

This inflates P@5 scores across all methods. However, it does
\textbf{not} affect MRR, which measures the rank of the \textbf{first}
relevant result---and the first result was consistently the correct
document for BM25. The inflated P@5 primarily narrows the apparent gap
between BM25 (26.7\%) and semantic (20.0\%), meaning the \textbf{true
performance gap is likely larger} than our reported numbers suggest.

For future work, we recommend using \textbf{document-level relevance
labels} (each query annotated with specific document IDs that contain
the answer) rather than keyword-level matching, and reporting
\textbf{nDCG} alongside P@K and MRR.

\subsection{8. Limitations}\label{limitations}

\begin{itemize}
\tightlist
\item
  \textbf{Dataset scale}: 41 documents and 105 queries provide
  sufficient power for statistical significance (Wilcoxon
  p\textless0.001), but larger corpora (\textgreater200 docs) would
  enable more granular ablation studies.
\item
  \textbf{Single embedding model}: Only all-MiniLM-L6-v2 was tested.
  Domain-specific models (e.g., BGE-Med, PubMedBERT) may perform
  differently.
\item
  \textbf{Chinese + Legal only}: Validation on additional languages and
  domains (e.g., finance, biology) would strengthen the
  language-mismatch hypothesis.
\end{itemize}

\subsection{9. Conclusion}\label{conclusion}

This paper presents the first controlled 2×2 crossover experiment
isolating language mismatch from domain specialization in non-English
RAG retrieval. Our key finding is that \textbf{the language of the
embedding model---not the retrieval algorithm, not the domain---is the
single largest controllable factor in RAG retrieval quality}, accounting
for a 13-20 percentage point penalty when mismatched. With a
language-appropriate embedding model, hybrid retrieval achieves
P@3=75.65\%, outperforming BM25 alone by 22 percentage points.

We release Medical-RAG-Bench as a resource for continued research, and
encourage the RAG community to treat embedding model language as a
first-class design consideration, not an afterthought.

\subsection{References}\label{references}

\begin{enumerate}
\def\labelenumi{\arabic{enumi}.}
\tightlist
\item
  Robertson, S., \& Zaragoza, H. (2009). The probabilistic relevance
  framework: BM25 and beyond.
\item
  Cormack, G. V., et al.~(2009). Reciprocal rank fusion outperforms
  condorcet and individual rank learning methods.
\item
  Lewis, P., et al.~(2020). Retrieval-Augmented Generation for
  Knowledge-Intensive NLP Tasks.
\item
  Reimers, N., \& Gurevych, I. (2019). Sentence-BERT: Sentence
  Embeddings using Siamese BERT-Networks.
\end{enumerate}

\end{document}
